%% Overhead foils starter — the standard LaTeX `slides' class.
%%
%% slides is part of core LaTeX and is NOT included in this download: it ships
%% with every modern TeX distribution (and with latextemplates.com's
%% compiler), so bundling a copy would only risk pinning you to an older
%% version than your own system already has.
%%
%% `slides' is the LaTeX2e successor to SliTeX, and produces large-format
%% overhead transparencies: oversized sans-serif type, one idea per foil. It
%% is the right class for a plain, high-contrast talk with no theme
%% machinery — use beamer instead if you want overlays and navigation.
%%
%% Class documentation: texdoc slides   |   https://ctan.org/pkg/slides
%% Licence: slides is LPPL-1.3c, (c) The LaTeX Project.
\documentclass[landscape]{slides}

\usepackage[T1]{fontenc}

%% slides' `landscape' option swaps \paperwidth and \paperheight, but
%% pdflatex sizes its output from \pdfpagewidth/\pdfpageheight, which the
%% class does not touch. Without these two lines the foils are typeset
%% landscape onto a portrait page, leaving the content in the top half.
\pdfpagewidth=\paperwidth
\pdfpageheight=\paperheight

\begin{document}

%% Each `slide' environment is one foil. Keep them short: the whole point of
%% the large type is that it stays readable from the back of a room.
\begin{slide}
\begin{center}
{\Large Your Talk Title Here}

\vspace{2em}
Ada Lovelace

\vspace{1em}
Institute of Analytical Engines

\vspace{1em}
\today
\end{center}
\end{slide}

\begin{slide}
\begin{center}\Large Outline\end{center}

\begin{itemize}
\item The problem
\item What we did
\item What we found
\item What is next
\end{itemize}
\end{slide}

\begin{slide}
\begin{center}\Large The Problem\end{center}

State it in one sentence.

\begin{itemize}
\item Why it matters
\item Why it is hard
\end{itemize}
\end{slide}

\begin{slide}
\begin{center}\Large Results\end{center}

\begin{center}
\begin{tabular}{lr}
Baseline & 0.71\\
Ours     & 0.86\\
\end{tabular}
\end{center}

One number the audience should remember.
\end{slide}

%% An `overlay' prints on top of the preceding slide — the classic
%% transparency trick of revealing an extra layer.
\begin{overlay}
\begin{center}\Large \ldots and it generalises.\end{center}
\end{overlay}

\begin{slide}
\begin{center}\Large Conclusion\end{center}

\begin{itemize}
\item What we showed
\item What comes next
\end{itemize}

\vspace{2em}
\begin{center}Questions?\end{center}
\end{slide}

\end{document}
